\section{Named exact covers and tomographic fibers}\label{sec:model}

\subsection{Finite placement model}

Let \(T\subset\ZZ^2\) be a finite target cell set.  A library consists of
nonempty, named, finite cell-list pieces.  Legal poses are obtained from a
supplied finite orientation group \(H\), with \(|H|\leq h\), followed by
integer translations; every legal pose must be contained in \(T\).
Explicitly, each \(g\in H\) acts by a lattice automorphism
\(g:\ZZ^2\to\ZZ^2\), the input supplies the action (or names one fixed
action), and the pose catalogue consists of the distinct cell masks
\(gP+t\subseteq T\).  Repeated masks produced by symmetries are identified.
Names remain distinct even if two pieces have the same shape or area.

Fix a selected named set
\[
  E=\{P_0,\ldots,P_r\},\qquad \sum_{i=0}^r |P_i|=|T|.
\]
The area equality is part of the decision domain.  Area-incompatible
selections are rejected during preprocessing.  Write \(\PP_i\) for the
finite catalogue of legal poses of \(P_i\).

For a cell set \(A\subseteq T\), let \(X(A)\) be its joint row-and-column
margin vector.  A one-copy placement tuple is
\[
  x=(A_0,\ldots,A_r)\in\prod_{i=0}^r\PP_i.
\]
Its cell coverage is
\[
  c_x(u)=|\{i:u\in A_i\}|,\qquad u\in T.
\]

\begin{definition}[Margin fiber and exact-cover sublocus]\label{def:fiber}
For the target margin \(X(T)\), define
\[
  \FFib_E=\left\{x:\sum_i X(A_i)=X(T)\right\}
\]
and
\[
  \TT_E=\{x\in\FFib_E:c_x=\one_T\}.
\]
Thus \(\TT_E\subseteq\FFib_E\).  Points of \(\FFib_E\setminus\TT_E\)
may have overlaps and uncovered cells, but every chosen pose still lies in
the target.
\end{definition}

The row and column margins are aggregate: the contribution of a piece name
is not separately observed.  This differs from typed projections that record
one margin vector per tile type.  It also differs from exact tiling models in
which disjointness is imposed before the projection constraints.

\subsection{The replacement graph}

\begin{definition}[Exact-two replacement graph]\label{def:graph}
The graph \(\GG_E^{(2)}\) has vertex set \(\FFib_E\).  Two distinct tuples
are adjacent precisely when they differ in exactly two named coordinates.
\end{definition}

An edge is one graph step, although its endpoints have Hamming distance two
as named placement tuples.  We also use the explicitly labelled variant
\(\GG_E^{(\leq2)}\), which adds edges between tuples differing in one named
coordinate.  These graphs need not agree.  Every unqualified repair
statement in this paper refers to \(\GG_E^{(2)}\).

\begin{definition}[Three exactness predicates]\label{def:predicates}
For every area-compatible selected set \(E\), define the following three
Boolean predicates:
\begin{itemize}
  \item \emph{support-exact} if
        \(\FFib_E\neq\varnothing\) exactly when
        \(\TT_E\neq\varnothing\);
  \item \emph{fiber-pure} if \(\FFib_E=\TT_E\);
  \item \emph{repair-exact} if every connected component of
        \(\GG_E^{(2)}\) meets \(\TT_E\).
\end{itemize}
If repair-exactness holds, the repair radius is
\[
  \rho(E)=\max_{x\in\FFib_E}\dist_{\GG_E^{(2)}}(x,\TT_E).
\]
\end{definition}

Equivalently, a finite library has any one of these properties on a declared
selection domain \(\mathcal D\) when the corresponding predicate holds for
every \(E\in\mathcal D\).  Finite-library tables in this paper always state
their domain explicitly.  Empty fibers satisfy all three predicates
vacuously and have radius zero.
If a component is disjoint from \(\TT_E\), its vertices are called
\emph{trapped}; they are not assigned a finite repair distance.

\subsection{Input-size contract}

The target and piece masks are explicit finite cell lists.  The orientation
group is supplied explicitly by its lattice action, or by a fixed named
lattice action of size at most \(h\).
Translations and margin vectors use binary-encoded integer coordinates.
The input length \(|I|\) includes these lists and coordinates.  The theorem
does not address arbitrary continuous bodies, oracle-described shapes,
unlimited tile multiplicities, or anonymous tile types.
